\section{The intrinsic polarization and a projective moduli completion}
\label{sec:polarization-v154}
Put $P=\Gr(2,\Sym^2V)$, $N=\binom{n+1}{2}$,
$M=\binom N2$, and $D=n^2+2n-4$. Fix temporarily an auxiliary
$n$-space $U$ and put $E=\Hom(U,V)^*$. Write $L=\OO_P(1)$ for
the Pluecker line and $\Gamma=\Phi(P)\subset\Gr(M,\Sym^D E)$ for
the closed image of Theorem~\ref{thm:finite-parameter-immersion}.
The group in this section is $\operatorname{SL}(V)$, acting by
congruence on pencils and by its induced action on $E$, with $U$
fixed. It is not the full group $\operatorname{GL}(E)$.

\begin{theorem}[Polarized completion of the effective failure moduli]
\label{thm:polarized-completion-v154}
There is an $\operatorname{SL}(V)$-equivariant isomorphism
\begin{equation}\label{eq:polarization-v154}
 \Phi^*\OO_{\Gr(M,\Sym^D E)}(1)\simeq L^{\otimes M}
                 \otimes C_U,
\end{equation}
where $C_U$ is a constant one-dimensional space on which
$\operatorname{SL}(V)$ acts trivially. Consequently the stable and
semistable loci of $\Gamma$, for the restricted Pluecker
linearization, are respectively the images of $P^s$ and $P^{ss}$.
The quotient defined by the complete restricted section ring is
canonically
\begin{equation}\label{eq:git-completion-v154}
 \overline{\mathcal M}^{\mathrm{fail}}_n
 =\Proj\!\left(\bigoplus_{q\geq0}
                  H^0(P,L^{\otimes q})^{\operatorname{SL}(V)}\right)
 \simeq\Gamma/\!/\operatorname{SL}(V).
\end{equation}
It is a normal projective variety. It contains the geometric quotient
of the simple-discriminant pencils as a dense open subvariety, and
has dimension $n-3$. The complement of that open is the image of the
semistable discriminant-zero locus.

The intrinsic inverse identifies the effective semistable failure
stack with $[P^{ss}/\PGL(V)]$ and gives a moduli morphism to
\eqref{eq:git-completion-v154}. On framed charts the failure-algebra
family extends over all $P$, remains finite locally free, and
commutes with arbitrary base change. These assertions do not require,
and do not assert, the existence of a universal ordinary algebra on
the coarse quotient or properness of the entire quotient stack.
\end{theorem}
\begin{proof}
Let $\mathcal K$ be the tautological relation bundle pulled back by
$\Phi$. The coefficient factorization
\eqref{eq:finite-coefficient-factorization} identifies it with
$\mathcal B\otimes F(U)$, where $\mathcal B$ is the quotient-volume
coefficient line. Exterior duality gives
$\mathcal B\simeq L^{-1}$ up to a fixed determinant character of
$\operatorname{GL}(V)$. That character is trivial on
$\operatorname{SL}(V)$. Taking the dual determinant gives
\eqref{eq:polarization-v154}, with the harmless constant factor
$(\det F(U))^*$ and the corresponding fixed determinant line of $V$.
A trivialization of the constant factor identifies the section rings;
a different trivialization rescales each graded piece and has no
effect on its projective spectrum. The section ring for $L^M$ is
the $M$th Veronese of the section ring for $L$. Taking invariants
commutes with this operation. The Hilbert--Mumford weights are
multiplied by $M>0$, so stability and semistability are unchanged.
This also explains why one uses the complete section ring on
$\Gamma$, rather than assuming surjectivity of unspecified ambient
invariant restrictions.

Finite generation and the projective good quotient are the reductive
GIT construction of \cite[Chapter~1]{MFK1994}. The complete section
ring of the normal projective variety $P$ is integrally closed: it
is the intersection of the valuation conditions for rational sections
of powers of $L$ at the prime divisors of $P$. Its invariant subring
is integrally closed as well. Indeed an element of its fraction
field integral over the invariants is integral over the original
ring, belongs to that ring by normality, and is invariant. This proves
normality of the projective quotient.

For completeness, the simple-discriminant open is nonempty and
stable. In a frame $(Q_0,Q_1)$ of a pencil, take the discriminant of
$\det(sQ_0+tQ_1)$. Under a frame change $h\in\operatorname{GL}_2$
it has weight $(\det h)^{n(n-1)}$, and it is invariant under
$\operatorname{SL}(V)$. Thus it is a nonzero invariant section of
$L^{n(n-1)}$. It is nonzero at a diagonal pencil with $n$ distinct
roots, so its nonvanishing locus is semistable and saturated.
A pencil there has a nonsingular member. Relative to that member,
the other member is self-adjoint with distinct eigenvalues;
orthogonal diagonalization identifies its congruence class with an
unordered configuration of $n$ distinct points of $\Pj^1$.
The latter configurations are stable for the usual binary-form
linearization: the root-multiplicity criterion is strict since
$1<n/2$ for $n\geq3$. Their quotient is geometric and separated.
It follows that the congruence orbits are closed in the saturated
nonvanishing locus. Their stabilizers are finite: the induced
projective transformation permutes a set of at least three points,
and the kernel consists of diagonal sign changes modulo scalars.
They are therefore stable. This gives a geometric dense open of
dimension $n-3$. The complement is the image of the invariant
section's zero locus by the good-quotient property. This reasoning
uses classical pencil classification and GIT, not a new classification
of binary forms; compare \cite{FevolaMandelshtamSturmfels,MFK1994}.

Finally apply Theorem~\ref{thm:pencil-stack-equivalence} and the
rigidification results of Section~\ref{sec:rigidification-v151}.
The Pluecker polarization is taken on the recovered pencil, not on
an arbitrary first-relation orbit. A sufficiently divisible power
of $L$, for example $L^n$, has trivial central character and descends
from $\operatorname{SL}(V)$ to $\PGL(V)$. A change of auxiliary frame
changes the constant factor in \eqref{eq:polarization-v154}, not the
resulting projective quotient. On the framed Grassmannian the extended
algebras are exactly those of Corollary~\ref{cor:finite-flat-family}.
Their equivariance yields the stated classifying morphisms. Removing
ineffective inertia does not turn these into a universal ordinary
algebra on the coarse space, and no such descent is used.
\end{proof}

\begin{remark}[What is compactified]\label{rem:git-scope-v154}
The completion is the classical polarized pencil quotient recovered
from the intrinsic failure invariant, with its polarization verified
by \eqref{eq:polarization-v154}. It is not a new construction of
classical GIT. Its role is to supply a proper coarse target for the
effective invariant and a specified class of semistable limiting
families. Strictly semistable points represent closed-orbit
representatives, not every orbit in their fibres. In particular Segre
data on the stack need not descend to an invariant separating all
points of a coarse-quotient fibre. For $n=3$ the quotient is a point
and its discriminant boundary can be empty. The full
$\operatorname{GL}(E)$ common limit of
Theorem~\ref{thm:universal-boundary-v152} is a different construction.
\end{remark}

\section{Spectral divisor fibres and the recovery of degeneration type}
\label{sec:spectral-divisor-v154}
The preceding completion concerns the native congruence action. To
retain degeneration type before passage to a coarse quotient, we
associate a tower of binary divisor incidences to a regular pencil.
This construction is intrinsic after the unmarked inverse and is
unchanged by a simultaneous change of the pencil and source frames.

Let $\Lambda=\Pj(R)$ be the pencil line and let
$\mathcal T_R$ be the cokernel of the universal symmetric map on
$\Lambda$. A regular pencil means that its determinant is not
identically zero. Thus $\mathcal T_R$ is a torsion sheaf. Define
$D_j(R)$ by
\begin{equation}\label{eq:spectral-Fitting-v154}
 \II_{D_j(R)}=\Fitt_j(\mathcal T_R),\qquad 0\leq j\leq n,
\end{equation}
where $D_n=\varnothing$. These are effective divisors on the smooth
curve $\Lambda$. Write $s_j=\deg D_j$. If $s_j>0$, let $G_j$ be
its binary equation and put
\[
 J_j=G_j H^0(\Lambda,\OO_\Lambda(1))
       \subset H^0(\Lambda,\OO_\Lambda(s_j+1)).
\]
Changing the scalar representative of $G_j$ does not change $J_j$.
Let $\mathcal F_j(R)$ be the whole scheme fibre at $J_j$ of the
degree-one divisor normalization $\nu_1$, with $r=2$ and $D=s_j+1$.
For $s_j=0$ put $\mathcal F_j=\varnothing$.

\begin{theorem}[A fibre-theoretic Segre invariant]\label{thm:spectral-fibres-v154}
There is a natural identification of finite schemes over $\Lambda$
\begin{equation}\label{eq:fibre-divisor-v154}
                        \mathcal F_j(R)\simeq D_j(R).
\end{equation}
At a spectral point $x$, write the positive elementary divisors of
$\mathcal T_R$ as
$\lambda_{x,1}\geq\cdots\geq\lambda_{x,r_x}>0$, and extend this
sequence by zeroes. Then
\begin{equation}\label{eq:fibre-length-v154}
 (\mathcal F_j(R))_x=\Spec\C[\epsilon]/(\epsilon^{b_{x,j}}),
 \qquad b_{x,j}=\sum_{i>j}\lambda_{x,i},
\end{equation}
with the empty-scheme convention when $b_{x,j}=0$. Consequently
\begin{equation}\label{eq:Segre-difference-v154}
                 \lambda_{x,j+1}=b_{x,j}-b_{x,j+1}.
\end{equation}
The tower, with its maps to the same intrinsic spectral line,
therefore determines the full Segre symbol. Conversely, the
elementary-divisor partitions together with their spectral support
on that line determine the tower. The tower retains the projective
configuration of the support as well as the Segre symbol.
It is an invariant of the unmarked finite failure algebra.

The construction is compatible with base change on every family
stratum where the Fitting ideals in \eqref{eq:spectral-Fitting-v154}
define relative effective Cartier divisors of fixed degrees. On such
a stratum $\mathcal F_j$ is finite locally free of degree $s_j$.
These hypotheses are part of the relative statement; no assertion of
flatness across a jump of $s_j$ is intended.
\end{theorem}
\begin{proof}
Over the discrete valuation ring $\OO_{\Lambda,x}$, Smith normal
form gives diagonal entries with exponents
$\lambda_{x,r_x},\ldots,\lambda_{x,1}$ and $n-r_x$ zero
exponents. The minors of size $n-j$ have greatest common divisor
of order $\sum_{i>j}\lambda_{x,i}$. This proves the formula for
$D_j$, including its nonreduced structure. It is the usual
Fitting-ideal description of elementary divisors, used here on the
intrinsic spectral line.

The residual two-space $H^0(\OO_\Lambda(1))$ is gcd-free. Hence
Theorem~\ref{thm:binary-fibres-v153} identifies the degree-one
normalization fibre at $J_j$ with the scheme of linear factors of
$G_j$. On an affine chart write a monic linear factor as $z-a$.
Monic division gives $G_j=(z-a)H+G_j(a)$, so the entire factor
scheme is $G_j(a)=0$, not its reduction. These descriptions agree
on overlaps of the projective line. They give
\eqref{eq:fibre-divisor-v154} and the local algebra
\eqref{eq:fibre-length-v154}. Taking successive differences proves
\eqref{eq:Segre-difference-v154}. The common map to $\Lambda$ is
essential: forgetting it separately in each member of the tower
would discard the alignment of the spectral supports.

Pencil frame changes act on $\Lambda$ by projective transformations;
source congruence changes give isomorphic cokernels and preserve all
Fitting ideals. The intrinsic inverse of
Theorem~\ref{thm:artin-local-inverse-v146} therefore defines this
invariant from the unmarked algebra. The equivalence with Segre data
uses the classical elementary-divisor convention of
\cite[\S5]{FevolaMandelshtamSturmfels}.

For a family with the stated Cartier and degree conditions, the same
monic division argument works over any base algebra, including one
with nilpotents. It represents the divisor's functor of points, so
it commutes with base change. A relative effective divisor of fixed
degree on a projective-line bundle is finite locally free of that
degree. Local choices of $G_j$ differ by units and glue. Formation
of Fitting ideals itself commutes with base change; requiring them
to remain flat Cartier divisors is precisely what permits the
additional finite-flat conclusion.
\end{proof}

\begin{example}[Equal discriminant, different fibre towers]
\label{ex:segre-fibres-v154}
At a root of multiplicity four, the partitions $(3,1)$ and $(2,2)$
have the same zeroth fibre $\C[\epsilon]/(\epsilon^4)$.
Their first fibres are respectively $\C$ and
$\C[\epsilon]/(\epsilon^2)$; all higher fibres are empty there.
Thus the discriminant alone does not detect the distinction, whereas
the divisor-fibre tower does. Both are realized by regular symmetric
pencils: for a block of size $l$, take the reversal matrix $H_l$ and
the symmetric pair $(H_l,H_lJ_l(\alpha))$, where $J_l(\alpha)$ is
a Jordan block. Direct sums realize every indicated partition.
Additional simple blocks may be appended. The example is a statement
on the regular-pencil stack, not a claim that distinct Segre strata
always remain distinct after strictly semistable GIT identification.
\end{example}

\begin{remark}[The two divisor constructions]\label{rem:two-divisors-v154}
The primitive first relation of an actual pencil has exact common
factor $(\det T)^{n-1}$, by Lemma~\ref{lem:primitive-pencil}.
Its common-factor degree does not jump when the pencil changes
Segre type. The divisors $D_j$ above instead lie on the recovered
spectral line and define auxiliary binary first-relation subspaces
$J_j$. Thus \eqref{eq:Segre-difference-v154} relates pencil data to
intrinsically induced divisor ramification, not to a nonexistent
change in the conductor of the original exact-gcd locus. The
identification supplies a moduli consequence of the all-binary-fibre
theorem while respecting this distinction.
\end{remark}
